\documentclass[11pt,a4paper]{article}

%========================
% PACKAGES
%========================
\usepackage[utf8]{inputenc}
\usepackage[T1]{fontenc}
\usepackage[french]{babel}
\usepackage[a4paper,margin=1.85cm]{geometry}
\usepackage{amsmath,amssymb}
\usepackage{lmodern}
\usepackage{xcolor}
\usepackage{fancyhdr}
\usepackage{lastpage}
\usepackage{enumitem}
\usepackage{array,tabularx}
\usepackage{tikz}
\usepackage[most]{tcolorbox}
\usetikzlibrary{arrows.meta}

%========================
% COULEURS
%========================
\definecolor{bleuprepa}{RGB}{20,55,110}
\definecolor{vertprepa}{RGB}{20,105,70}
\definecolor{rougeprepa}{RGB}{160,35,35}
\definecolor{orangeprepa}{RGB}{180,100,20}
\definecolor{grisclair}{RGB}{245,247,250}

%========================
% PAGE
%========================
\setlength{\headheight}{14pt}
\pagestyle{fancy}
\fancyhf{}
\fancyhead[L]{Troisième -- TD complet}
\fancyhead[C]{Données, probabilités, proportionnalité, grandeurs}
\fancyhead[R]{Page \thepage/\pageref{LastPage}}
\fancyfoot[C]{Pierre Garcia-Buscaylet -- Mathématiques}

\setlength{\parindent}{0pt}
\setlength{\parskip}{0.35em}
\setlist[enumerate]{label=\textbf{\alph*)},leftmargin=*,itemsep=0.2em}

%========================
% COMMANDES
%========================
\newcommand{\classe}[2]{\ensuremath{[#1\,;#2[}}
\newcommand{\etoiles}[1]{\textcolor{orangeprepa}{#1}}
\newcommand{\Prob}{\mathbb{P}}

\newcounter{exo}
\newcounter{corr}

\newtcolorbox{exobox}[3]{
  enhanced,
  breakable,
  colback=white,
  colframe=bleuprepa,
  boxrule=0.7pt,
  arc=1mm,
  title={\stepcounter{exo}\textbf{Exercice \theexo\ -- #1}\hfill \etoiles{#2}\quad \textsc{#3}},
  fonttitle=\bfseries
}

\newtcolorbox{corrbox}[1]{
  enhanced,
  breakable,
  colback=vertprepa!3!white,
  colframe=vertprepa,
  boxrule=0.7pt,
  arc=1mm,
  title={\stepcounter{corr}\textbf{Correction de l'exercice \thecorr\ -- #1}},
  fonttitle=\bfseries
}

\newtcolorbox{bilanbox}[1]{
  enhanced,
  breakable,
  colback=grisclair,
  colframe=black!60,
  boxrule=0.7pt,
  arc=1mm,
  title={\textbf{#1}},
  fonttitle=\bfseries
}

\begin{document}

\begin{center}
{\Huge \textbf{TD complet de mathématiques}}\\[0.2cm]
{\Large \textbf{Classe de troisième}}\\[0.15cm]
{\Large Données -- Probabilités -- Proportionnalité -- Grandeurs mesurables}\\[0.2cm]
\rule{0.92\linewidth}{0.6pt}\\[0.15cm]
{\itshape 50 exercices progressifs avec corrigé complet}\\[0.15cm]
\rule{0.92\linewidth}{0.6pt}
\end{center}

\begin{bilanbox}{Objectif du TD}
Ce TD entraîne les élèves à lire, représenter et traiter des données, calculer des probabilités simples ou à deux épreuves, utiliser la proportionnalité et les pourcentages d'évolution, calculer avec des grandeurs mesurables et vérifier les unités. Les exercices sont organisés par familles : application, méthode, justification, réflexion et exercices longs.
\end{bilanbox}

\tableofcontents
\newpage

%================================================
\section{Exercices d'application}

\begin{exobox}{Effectifs et fréquences dans une classe}{$\star$}{APP}
Dans une classe de $28$ élèves, on relève le moyen principal utilisé pour venir au collège.
\[
\begin{array}{|c|c|}
\hline
\text{Moyen de transport} & \text{Effectif} \\
\hline
\text{À pied} & 7 \\
\text{Bus} & 12 \\
\text{Voiture} & 6 \\
\text{Vélo} & 3 \\
\hline
\end{array}
\]
\begin{enumerate}
\item Vérifier que l'effectif total est bien $28$.
\item Calculer la fréquence des élèves venant en bus.
\item Calculer la fréquence des élèves venant à pied ou à vélo.
\item Écrire les deux fréquences précédentes en pourcentage.
\end{enumerate}
\end{exobox}

\begin{exobox}{Étendue d'une série brute}{$\star$}{APP}
On a relevé les températures maximales, en degrés Celsius, pendant dix jours :
\[
18,\ 21,\ 17,\ 23,\ 25,\ 20,\ 19,\ 24,\ 22,\ 16.
\]
\begin{enumerate}
\item Donner la valeur minimale de la série.
\item Donner la valeur maximale de la série.
\item Calculer l'étendue.
\item Interpréter cette étendue par une phrase précise.
\end{enumerate}
\end{exobox}

\begin{exobox}{Fréquences dans un tableau}{$\star$}{APP}
Une enquête porte sur la couleur préférée de $60$ élèves.
\[
\begin{array}{|c|c|}
\hline
\text{Couleur} & \text{Effectif} \\
\hline
\text{Bleu} & 18 \\
\text{Rouge} & 12 \\
\text{Vert} & 9 \\
\text{Noir} & 15 \\
\text{Autre} & 6 \\
\hline
\end{array}
\]
\begin{enumerate}
\item Calculer la fréquence de la modalité « Bleu ».
\item Calculer la fréquence de la modalité « Noir ».
\item Calculer la fréquence des élèves ayant répondu « Rouge » ou « Vert ».
\item Vérifier que la somme des fréquences vaut $1$.
\end{enumerate}
\end{exobox}

\begin{exobox}{Histogramme à classes de même amplitude}{$\star$}{APP}
On mesure la durée quotidienne de lecture de $100$ élèves.
\[
\begin{array}{|c|c|}
\hline
\text{Durée en minutes} & \text{Effectif} \\
\hline
\classe{0}{10} & 8 \\
\classe{10}{20} & 22 \\
\classe{20}{30} & 34 \\
\classe{30}{40} & 26 \\
\classe{40}{50} & 10 \\
\hline
\end{array}
\]
\begin{enumerate}
\item Vérifier l'effectif total.
\item Donner l'amplitude commune des classes.
\item Expliquer pourquoi un histogramme est adapté.
\item Calculer la fréquence des élèves lisant entre $20$ et $40$ minutes par jour.
\end{enumerate}
\end{exobox}

\begin{exobox}{Diagramme circulaire}{$\star\star$}{APP}
Un club compte $80$ adhérents. La répartition par activité est donnée ci-dessous.
\[
\begin{array}{|c|c|}
\hline
\text{Activité} & \text{Effectif} \\
\hline
\text{Football} & 32 \\
\text{Basket} & 20 \\
\text{Tennis} & 16 \\
\text{Athlétisme} & 12 \\
\hline
\end{array}
\]
\begin{enumerate}
\item Calculer la fréquence de chaque activité.
\item Calculer l'angle du secteur correspondant à chaque activité dans un diagramme circulaire.
\item Vérifier que la somme des angles vaut $360^\circ$.
\end{enumerate}
\end{exobox}

\begin{exobox}{Probabilité avec un dé équilibré}{$\star$}{APP}
On lance un dé équilibré à six faces numérotées de $1$ à $6$.
\begin{enumerate}
\item Donner l'univers de l'expérience.
\item Calculer la probabilité d'obtenir $5$.
\item Calculer la probabilité d'obtenir un nombre pair.
\item Calculer la probabilité d'obtenir un nombre supérieur ou égal à $4$.
\item Calculer la probabilité de ne pas obtenir $1$.
\end{enumerate}
\end{exobox}

\begin{exobox}{Probabilité dans une urne}{$\star$}{APP}
Une urne contient $4$ boules rouges, $3$ boules bleues et $5$ boules vertes. On tire une boule au hasard.
\begin{enumerate}
\item Calculer le nombre total de boules.
\item Calculer la probabilité de tirer une boule rouge.
\item Calculer la probabilité de tirer une boule bleue ou verte.
\item Calculer la probabilité de ne pas tirer une boule verte.
\end{enumerate}
\end{exobox}

\begin{exobox}{Deux pièces équilibrées}{$\star\star$}{APP}
On lance deux pièces équilibrées. On note $P$ pour pile et $F$ pour face.
\begin{enumerate}
\item Lister toutes les issues possibles sous forme de couples.
\item Calculer la probabilité d'obtenir deux piles.
\item Calculer la probabilité d'obtenir exactement une pile.
\item Calculer la probabilité d'obtenir au moins une face.
\end{enumerate}
\end{exobox}

\begin{exobox}{Fonction linéaire et vitesse constante}{$\star$}{APP}
Un mobile avance à vitesse constante de $5$ mètres par seconde. On note $x$ le temps en secondes et $d(x)$ la distance parcourue en mètres.
\begin{enumerate}
\item Écrire la fonction linéaire qui modélise la situation.
\item Calculer $d(8)$.
\item Déterminer la distance parcourue en $1$ minute.
\item Déterminer le temps nécessaire pour parcourir $250$ mètres.
\end{enumerate}
\end{exobox}

\begin{exobox}{Pourcentages d'évolution}{$\star\star$}{APP}
Un article coûte $120$ euros.
\begin{enumerate}
\item Calculer son prix après une augmentation de $15\%$.
\item Calculer son prix après une diminution de $20\%$.
\item Donner le coefficient multiplicateur d'une hausse de $8\%$.
\item Donner le coefficient multiplicateur d'une baisse de $35\%$.
\end{enumerate}
\end{exobox}

\begin{exobox}{Volume d'une boule}{$\star$}{APP}
On considère une boule de rayon $6$ cm.
\begin{enumerate}
\item Rappeler la formule du volume d'une boule de rayon $r$.
\item Calculer son volume exact.
\item Donner une valeur approchée au dixième de cm$^3$ près.
\item Vérifier que l'unité finale est cohérente.
\end{enumerate}
\end{exobox}

\begin{exobox}{Vitesse et conversion}{$\star\star$}{APP}
Une voiture roule à $72$ km/h.
\begin{enumerate}
\item Convertir cette vitesse en m/s.
\item Calculer la distance parcourue en $3$ secondes.
\item Calculer la distance parcourue en $12$ secondes.
\item Expliquer pourquoi il fallait convertir la vitesse avant de calculer ces distances.
\end{enumerate}
\end{exobox}

%================================================
\section{Exercices de méthode}

\begin{exobox}{Retrouver des effectifs à partir de fréquences}{$\star\star$}{METH}
Dans un collège de $480$ élèves, une enquête donne les fréquences suivantes pour le trajet du matin.
\[
\begin{array}{|c|c|}
\hline
\text{Moyen} & \text{Fréquence} \\
\hline
\text{Bus} & 0{,}45 \\
\text{Voiture} & 0{,}25 \\
\text{À pied} & 0{,}20 \\
\text{Vélo} & 0{,}10 \\
\hline
\end{array}
\]
\begin{enumerate}
\item Vérifier que les fréquences forment une série complète.
\item Calculer l'effectif correspondant à chaque moyen de transport.
\item Contrôler que la somme des effectifs obtenus vaut $480$.
\item Présenter les résultats dans un tableau d'effectifs.
\end{enumerate}
\end{exobox}

\begin{exobox}{Lire un diagramme en bâtons décrit par un tableau}{$\star\star$}{METH}
Le tableau suivant correspond à un diagramme en bâtons donnant les pointures de $36$ élèves.
\[
\begin{array}{|c|c|}
\hline
\text{Pointure} & \text{Effectif} \\
\hline
37 & 2 \\
38 & 5 \\
39 & 7 \\
40 & 8 \\
41 & 6 \\
42 & 5 \\
43 & 3 \\
\hline
\end{array}
\]
\begin{enumerate}
\item Calculer le nombre d'élèves chaussant au moins du $40$.
\item Calculer le nombre d'élèves chaussant au plus du $42$.
\item Calculer le nombre d'élèves chaussant entre $38$ et $41$, bornes incluses.
\item Calculer la fréquence des élèves chaussant au moins du $41$.
\end{enumerate}
\end{exobox}

\begin{exobox}{Construire un histogramme complet}{$\star\star$}{METH}
On regroupe les masses, en kilogrammes, de $50$ colis dans le tableau suivant.
\[
\begin{array}{|c|c|}
\hline
\text{Masse en kg} & \text{Effectif} \\
\hline
\classe{0}{5} & 6 \\
\classe{5}{10} & 14 \\
\classe{10}{15} & 18 \\
\classe{15}{20} & 9 \\
\classe{20}{25} & 3 \\
\hline
\end{array}
\]
\begin{enumerate}
\item Vérifier que les classes ont même amplitude.
\item Donner l'échelle possible pour tracer l'axe vertical.
\item Construire l'histogramme sur papier.
\item Calculer la fréquence des colis de masse au moins $10$ kg.
\end{enumerate}
\end{exobox}

\begin{exobox}{Étendue à partir d'un histogramme}{$\star\star$}{METH}
Un histogramme de tailles regroupe des élèves dans les classes suivantes :
\[
\classe{140}{150},\quad \classe{150}{160},\quad \classe{160}{170},\quad \classe{170}{180},\quad \classe{180}{190}.
\]
Les effectifs correspondants sont $3$, $9$, $15$, $10$, $3$.
\begin{enumerate}
\item Construire le tableau de la série.
\item Calculer l'effectif total.
\item Donner une estimation de l'étendue à partir des bornes des classes.
\item Expliquer pourquoi cette étendue est une estimation et non forcément l'étendue exacte des données brutes.
\end{enumerate}
\end{exobox}

\begin{exobox}{Deux tirages avec remise}{$\star\star$}{METH}
Une urne contient une boule bleue $B$ et deux boules violettes $V_1$ et $V_2$. On tire deux fois avec remise.
\begin{enumerate}
\item Construire un tableau à double entrée des issues possibles.
\item Compter le nombre total d'issues.
\item Identifier les issues correspondant à deux boules violettes.
\item En déduire la probabilité de tirer deux boules violettes.
\end{enumerate}
\end{exobox}

\begin{exobox}{Deux enfants}{$\star\star$}{METH}
On suppose que, pour chaque naissance, la probabilité d'avoir une fille ou un garçon est la même. Un couple a deux enfants.
\begin{enumerate}
\item Lister toutes les issues possibles en tenant compte de l'ordre des naissances.
\item Calculer la probabilité d'avoir deux garçons.
\item Calculer directement la probabilité d'avoir au moins une fille.
\item Retrouver ce résultat avec l'événement contraire.
\end{enumerate}
\end{exobox}

\begin{exobox}{Simulation d'un dé}{$\star\star$}{METH}
On lance un dé équilibré $10\,000$ fois. Les effectifs obtenus sont les suivants.
\[
\begin{array}{|c|c|}
\hline
\text{Face} & \text{Effectif observé} \\
\hline
1 & 1650 \\
2 & 1712 \\
3 & 1640 \\
4 & 1688 \\
5 & 1661 \\
6 & 1649 \\
\hline
\end{array}
\]
\begin{enumerate}
\item Calculer la fréquence observée de chaque face au millième près.
\item Donner la probabilité théorique d'obtenir une face donnée.
\item Comparer les fréquences observées à la probabilité théorique.
\item Expliquer en une phrase le phénomène de stabilisation des fréquences.
\end{enumerate}
\end{exobox}

\begin{exobox}{Coefficient multiplicateur global}{$\star\star\star$}{METH}
Un prix subit successivement une hausse de $10\%$, puis une baisse de $10\%$.
\begin{enumerate}
\item Donner le coefficient multiplicateur de la hausse.
\item Donner le coefficient multiplicateur de la baisse.
\item Calculer le coefficient multiplicateur global.
\item Dire si le prix final est égal au prix initial, puis justifier.
\end{enumerate}
\end{exobox}

\begin{exobox}{Fonction linéaire et tableau}{$\star\star$}{METH}
Une grandeur $y$ est proportionnelle à une grandeur $x$. On sait que $y=18$ lorsque $x=6$.
\begin{enumerate}
\item Calculer le coefficient de proportionnalité.
\item Écrire la fonction linéaire associée.
\item Compléter le tableau :
\[
\begin{array}{|c|c|c|c|c|}
\hline
x & 2 & 5 & 9 & 12 \\
\hline
y & & & & \\
\hline
\end{array}
\]
\item Déterminer $x$ lorsque $y=45$.
\end{enumerate}
\end{exobox}

\begin{exobox}{Thalès et proportionnalité}{$\star\star\star$}{METH}
Dans un triangle $ABC$, les points $M$ et $N$ appartiennent respectivement aux segments $[AB]$ et $[AC]$, et les droites $(MN)$ et $(BC)$ sont parallèles. On donne $AM=4$ cm, $AB=10$ cm et $AC=15$ cm.
\begin{enumerate}
\item Écrire l'égalité de Thalès utile.
\item Calculer le rapport $\dfrac{AM}{AB}$.
\item Calculer $AN$.
\item Interpréter le résultat comme une réduction du triangle $ABC$.
\end{enumerate}
\end{exobox}

\begin{exobox}{Débit et unités}{$\star\star$}{METH}
Un robinet a un débit constant de $18$ L/min.
\begin{enumerate}
\item Calculer le volume écoulé en $7$ minutes.
\item Convertir ce volume en m$^3$.
\item Calculer le temps nécessaire pour remplir une cuve de $540$ L.
\item Vérifier la cohérence des unités dans chaque calcul.
\end{enumerate}
\end{exobox}

\begin{exobox}{Assemblage de solides}{$\star\star\star$}{METH}
Un solide est composé d'un cylindre de rayon $3$ cm et de hauteur $12$ cm, surmonté d'une demi-boule de même rayon.
\begin{enumerate}
\item Calculer le volume du cylindre.
\item Calculer le volume de la demi-boule.
\item Calculer le volume total exact.
\item Donner une valeur approchée au cm$^3$ près.
\end{enumerate}
\end{exobox}

%================================================
\section{Exercices de justification et de démonstration}

\begin{exobox}{Justifier une formule de fréquence}{$\star\star$}{DEM}
Dans une série statistique d'effectif total $N$, une valeur possède un effectif $n$.
\begin{enumerate}
\item Expliquer pourquoi la fréquence associée est $f=\dfrac{n}{N}$.
\item Justifier que $0\leq f\leq 1$.
\item Expliquer pourquoi la somme des fréquences de toutes les modalités vaut $1$.
\item Donner une interprétation lorsque $f=0{,}25$.
\end{enumerate}
\end{exobox}

\begin{exobox}{Justifier l'étendue}{$\star\star$}{DEM}
On considère une série statistique de plus petite valeur $m$ et de plus grande valeur $M$.
\begin{enumerate}
\item Rappeler la définition de l'étendue.
\item Expliquer pourquoi l'étendue est toujours positive ou nulle.
\item Donner une situation où l'étendue vaut $0$.
\item Expliquer ce que mesure l'étendue et ce qu'elle ne mesure pas.
\end{enumerate}
\end{exobox}

\begin{exobox}{Justifier l'événement contraire}{$\star\star$}{DEM}
On considère un événement $A$ dans une expérience aléatoire finie en situation d'équiprobabilité.
\begin{enumerate}
\item Expliquer ce que signifie l'événement contraire $\overline A$.
\item Justifier que les issues favorables à $A$ et celles favorables à $\overline A$ ne se recouvrent pas.
\item Justifier que toutes les issues appartiennent soit à $A$, soit à $\overline A$.
\item En déduire que $\Prob(\overline A)=1-\Prob(A)$.
\end{enumerate}
\end{exobox}

\begin{exobox}{Justifier une situation de proportionnalité}{$\star\star$}{DEM}
On donne le tableau suivant.
\[
\begin{array}{|c|c|c|c|c|}
\hline
x & 3 & 5 & 8 & 12 \\
\hline
y & 7{,}5 & 12{,}5 & 20 & 30 \\
\hline
\end{array}
\]
\begin{enumerate}
\item Calculer chaque quotient $\dfrac{y}{x}$.
\item Conclure sur la proportionnalité.
\item Donner le coefficient de proportionnalité.
\item Écrire la fonction linéaire associée.
\end{enumerate}
\end{exobox}

\begin{exobox}{Prouver qu'une situation n'est pas proportionnelle}{$\star\star$}{DEM}
On donne le tableau suivant.
\[
\begin{array}{|c|c|c|c|}
\hline
x & 2 & 4 & 8 \\
\hline
y & 5 & 9 & 17 \\
\hline
\end{array}
\]
\begin{enumerate}
\item Calculer $\dfrac{5}{2}$, $\dfrac{9}{4}$ et $\dfrac{17}{8}$.
\item Comparer ces quotients.
\item Justifier rigoureusement que le tableau n'est pas proportionnel.
\item Expliquer pourquoi il ne suffit pas que $y$ augmente quand $x$ augmente.
\end{enumerate}
\end{exobox}

\begin{exobox}{Démontrer le coefficient d'une hausse}{$\star\star$}{DEM}
On considère une valeur initiale $V$.
\begin{enumerate}
\item Écrire la quantité ajoutée lors d'une hausse de $p\%$.
\item Écrire la valeur finale sous la forme $V+\dfrac{p}{100}V$.
\item Factoriser par $V$.
\item En déduire que le coefficient multiplicateur est $1+\dfrac{p}{100}$.
\end{enumerate}
\end{exobox}

\begin{exobox}{Démontrer le coefficient d'une baisse}{$\star\star$}{DEM}
On considère une valeur initiale $V$.
\begin{enumerate}
\item Écrire la quantité retirée lors d'une baisse de $p\%$.
\item Écrire la valeur finale sous la forme $V-\dfrac{p}{100}V$.
\item Factoriser par $V$.
\item En déduire que le coefficient multiplicateur est $1-\dfrac{p}{100}$.
\end{enumerate}
\end{exobox}

\begin{exobox}{Contrôle d'unités dans une vitesse}{$\star\star\star$}{DEM}
On rappelle que $v=\dfrac{d}{t}$.
\begin{enumerate}
\item Si $d$ est en mètres et $t$ en secondes, donner l'unité de $v$.
\item Si $v$ est en m/s et $t$ en secondes, expliquer pourquoi $d=v\times t$ est en mètres.
\item Convertir $1$ km/h en m/s en partant de $1$ km $=1000$ m et $1$ h $=3600$ s.
\item En déduire pourquoi on divise par $3{,}6$ pour passer de km/h à m/s.
\end{enumerate}
\end{exobox}

%================================================
\section{Exercices de réflexion}

\begin{exobox}{Histogramme et piège de lecture}{$\star\star\star$}{PREP}
Une enquête porte sur $120$ temps d'attente à un guichet.
\[
\begin{array}{|c|c|}
\hline
\text{Temps en minutes} & \text{Effectif} \\
\hline
\classe{0}{5} & 10 \\
\classe{5}{10} & 18 \\
\classe{10}{15} & 42 \\
\classe{15}{20} & 30 \\
\classe{20}{25} & 20 \\
\hline
\end{array}
\]
\begin{enumerate}
\item Quelle est la classe modale, c'est-à-dire la classe de plus grand effectif ?
\item Calculer la fréquence des temps d'attente inférieurs à $10$ minutes.
\item Calculer la fréquence des temps d'attente au moins égaux à $15$ minutes.
\item Peut-on connaître exactement le nombre de personnes ayant attendu exactement $12$ minutes ? Justifier.
\end{enumerate}
\end{exobox}

\begin{exobox}{Probabilités et dénombrement non ordonné}{$\star\star\star$}{PREP}
Un sac contient $2$ jetons rouges $R_1,R_2$ et $2$ jetons bleus $B_1,B_2$. On tire deux jetons successivement avec remise.
\begin{enumerate}
\item Construire l'ensemble des $16$ issues en distinguant les jetons.
\item Calculer la probabilité d'obtenir deux jetons rouges.
\item Calculer la probabilité d'obtenir deux jetons de même couleur.
\item Expliquer pourquoi il faut distinguer les jetons pour un dénombrement parfaitement clair.
\end{enumerate}
\end{exobox}

\begin{exobox}{Comparer deux simulations}{$\star\star\star$}{PREP}
Deux élèves simulent des lancers de dé. Le premier effectue $60$ lancers et obtient $14$ fois la face $6$. Le second effectue $6000$ lancers et obtient $1015$ fois la face $6$.
\begin{enumerate}
\item Calculer la fréquence d'apparition du $6$ dans chaque simulation.
\item Donner la probabilité théorique d'obtenir $6$.
\item Quelle simulation est la plus proche de la probabilité théorique ?
\item Expliquer pourquoi le second résultat est généralement plus fiable pour estimer la probabilité.
\end{enumerate}
\end{exobox}

\begin{exobox}{Hausse puis baisse}{$\star\star\star$}{PREP}
Un prix initial de $200$ euros augmente de $25\%$, puis diminue de $20\%$.
\begin{enumerate}
\item Calculer le prix après la hausse.
\item Calculer le prix final après la baisse.
\item Calculer le coefficient multiplicateur global.
\item Expliquer pourquoi une hausse de $25\%$ suivie d'une baisse de $20\%$ ramène ici au prix initial.
\end{enumerate}
\end{exobox}

\begin{exobox}{Baisse puis hausse inverse ?}{$\star\star\star$}{PREP}
Un prix baisse de $30\%$.
\begin{enumerate}
\item Donner le coefficient multiplicateur de la baisse.
\item On veut retrouver le prix initial par une hausse. Si le prix après baisse est $70$, quel était le prix initial ?
\item Quel pourcentage de hausse faut-il appliquer à $70$ pour revenir à $100$ ?
\item Expliquer pourquoi ce pourcentage n'est pas $30\%$.
\end{enumerate}
\end{exobox}

\begin{exobox}{Thalès avec retour à la longueur initiale}{$\star\star\star$}{PREP}
Dans une configuration de Thalès, on a $(MN)\parallel(BC)$, $AM=6$ cm, $MB=9$ cm, $AN=8$ cm.
\begin{enumerate}
\item Calculer $AB$.
\item Écrire le rapport de réduction $\dfrac{AM}{AB}$.
\item Utiliser Thalès pour calculer $AC$.
\item En déduire $NC$.
\end{enumerate}
\end{exobox}

\begin{exobox}{Boîte cylindrique et boules}{$\star\star\star\star$}{PREP}
Une boîte cylindrique de hauteur $30$ cm et de rayon $5$ cm contient trois boules identiques de rayon $5$ cm empilées verticalement.
\begin{enumerate}
\item Calculer le volume de la boîte cylindrique.
\item Calculer le volume total des trois boules.
\item Calculer le volume restant dans la boîte.
\item Donner une valeur approchée au cm$^3$ près et expliquer si le résultat est cohérent.
\end{enumerate}
\end{exobox}

\begin{exobox}{Débit important}{$\star\star\star$}{PREP}
Le débit moyen d'une rivière est $328$ m$^3$/s.
\begin{enumerate}
\item Convertir $3$ minutes en secondes.
\item Calculer le volume d'eau écoulé en $3$ minutes en m$^3$.
\item Convertir ce volume en litres.
\item Écrire le résultat en notation scientifique.
\end{enumerate}
\end{exobox}

\begin{exobox}{Choisir la bonne méthode}{$\star\star\star\star$}{PREP}
Un collège organise un sondage auprès de $300$ élèves. $42\%$ pratiquent un sport en club. Parmi eux, $\frac{2}{7}$ pratiquent un sport collectif.
\begin{enumerate}
\item Calculer le nombre d'élèves pratiquant un sport en club.
\item Calculer le nombre d'élèves pratiquant un sport collectif parmi ceux qui pratiquent en club.
\item Calculer la fréquence, parmi tous les élèves interrogés, des élèves pratiquant un sport collectif en club.
\item Expliquer quelles opérations relèvent de la proportionnalité.
\end{enumerate}
\end{exobox}

\begin{exobox}{Données regroupées et incertitude}{$\star\star\star\star$}{PREP}
On étudie les durées de connexion de $200$ utilisateurs.
\[
\begin{array}{|c|c|}
\hline
\text{Durée en minutes} & \text{Effectif} \\
\hline
\classe{0}{20} & 25 \\
\classe{20}{40} & 50 \\
\classe{40}{60} & 70 \\
\classe{60}{80} & 40 \\
\classe{80}{100} & 15 \\
\hline
\end{array}
\]
\begin{enumerate}
\item Calculer la fréquence de la classe \classe{40}{60}.
\item Donner une estimation de l'étendue.
\item Peut-on savoir combien d'utilisateurs se sont connectés exactement $50$ minutes ?
\item Peut-on savoir combien se sont connectés au moins $40$ minutes ? Justifier.
\end{enumerate}
\end{exobox}

%================================================
\section{Exercices longs}

\begin{exobox}{Enquête statistique complète}{$\star\star\star\star$}{ORAUX}
Une mairie interroge $2500$ personnes : « À quel âge avez-vous trouvé un emploi correspondant à votre qualification ? »
\[
\begin{array}{|c|c|}
\hline
\text{Âge} & \text{Effectif} \\
\hline
\classe{18}{22} & 100 \\
\classe{22}{26} & 200 \\
\classe{26}{30} & 400 \\
\classe{30}{34} & 1100 \\
\classe{34}{38} & 700 \\
\hline
\end{array}
\]
\begin{enumerate}
\item Vérifier l'effectif total.
\item Calculer la fréquence de chaque classe.
\item Représenter les résultats par un histogramme.
\item Déterminer la classe la plus représentée.
\item Donner une estimation de l'étendue.
\item Rédiger une interprétation globale des résultats en trois phrases.
\end{enumerate}
\end{exobox}

\begin{exobox}{Jeu à deux épreuves}{$\star\star\star\star$}{ORAUX}
Une urne contient $1$ boule bleue et $2$ boules violettes. On tire deux fois avec remise. On gagne si l'on obtient au moins une boule violette.
\begin{enumerate}
\item Nommer les boules violettes $V_1,V_2$ et construire un tableau à double entrée.
\item Calculer la probabilité d'obtenir deux boules violettes.
\item Calculer la probabilité d'obtenir aucune boule violette.
\item En déduire la probabilité de gagner.
\item Vérifier le résultat par un dénombrement direct des issues gagnantes.
\end{enumerate}
\end{exobox}

\begin{exobox}{Trajet, vitesse et temps de réaction}{$\star\star\star\star$}{ORAUX}
Un conducteur roule à $80$ km/h. Lorsqu'il voit un obstacle, il met $1$ seconde avant de commencer à freiner. La distance de freinage est ensuite de $38$ m.
\begin{enumerate}
\item Convertir $80$ km/h en m/s.
\item Calculer la distance parcourue pendant le temps de réaction.
\item Calculer la distance totale d'arrêt.
\item Si le conducteur roulait à $90$ km/h avec la même distance de freinage, quelle serait la nouvelle distance de réaction ?
\item Comparer les deux situations.
\end{enumerate}
\end{exobox}

\begin{exobox}{Réservoir composé}{$\star\star\star\star$}{ORAUX}
Un réservoir est formé d'un cylindre de rayon $4$ dm et de hauteur $10$ dm, surmonté d'une demi-boule de même rayon.
\begin{enumerate}
\item Calculer le volume exact du cylindre en dm$^3$.
\item Calculer le volume exact de la demi-boule en dm$^3$.
\item Calculer le volume total en dm$^3$.
\item Convertir ce volume en litres.
\item Donner une valeur approchée au litre près.
\end{enumerate}
\end{exobox}

\begin{exobox}{Sondage, fréquences et évolution}{$\star\star\star\star$}{ORAUX}
Une association compte $600$ adhérents. La première année, $35\%$ des adhérents ont moins de $18$ ans. L'année suivante, le nombre total d'adhérents augmente de $12\%$ et la fréquence des moins de $18$ ans devient $40\%$.
\begin{enumerate}
\item Calculer le nombre d'adhérents de moins de $18$ ans la première année.
\item Calculer le nombre total d'adhérents l'année suivante.
\item Calculer le nombre d'adhérents de moins de $18$ ans l'année suivante.
\item Calculer l'augmentation en effectif des moins de $18$ ans.
\item Interpréter les résultats.
\end{enumerate}
\end{exobox}

\begin{exobox}{Fonction linéaire et géométrie}{$\star\star\star\star\star$}{ORAUX}
Une figure est agrandie par homothétie de coefficient $1{,}8$.
\begin{enumerate}
\item Écrire la fonction linéaire donnant une longueur agrandie $L'$ en fonction d'une longueur initiale $L$.
\item Calculer les images des longueurs $4$ cm, $7{,}5$ cm et $12$ cm.
\item Une longueur agrandie mesure $45$ cm. Retrouver la longueur initiale.
\item Expliquer pourquoi les longueurs sont proportionnelles.
\item Dire ce qui arriverait aux volumes si tous les solides étaient agrandis par le même coefficient, sans effectuer de calcul hors programme.
\end{enumerate}
\end{exobox}

\begin{exobox}{Interprétation d'une expérience simulée}{$\star\star\star\star\star$}{ORAUX}
Une simulation informatique lance une pièce équilibrée. Trois séries sont obtenues :
\[
\begin{array}{|c|c|c|}
\hline
\text{Nombre de lancers} & \text{Nombre de piles} & \text{Fréquence de piles} \\
\hline
20 & 14 &  \\
200 & 107 &  \\
20000 & 10084 &  \\
\hline
\end{array}
\]
\begin{enumerate}
\item Compléter la colonne des fréquences.
\item Donner la probabilité théorique d'obtenir pile.
\item Comparer chaque fréquence à la probabilité théorique.
\item Expliquer pourquoi la troisième simulation est la plus convaincante.
\item Rédiger une conclusion sur la stabilisation des fréquences.
\end{enumerate}
\end{exobox}

\begin{exobox}{Problème bilan transversal}{$\star\star\star\star\star$}{ORAUX}
Une course solidaire regroupe $800$ participants. Les âges sont regroupés ainsi :
\[
\begin{array}{|c|c|}
\hline
\text{Âge} & \text{Effectif} \\
\hline
\classe{10}{20} & 120 \\
\classe{20}{30} & 260 \\
\classe{30}{40} & 220 \\
\classe{40}{50} & 140 \\
\classe{50}{60} & 60 \\
\hline
\end{array}
\]
Chaque participant parcourt $5$ km. On estime qu'un participant met en moyenne $36$ minutes.
\begin{enumerate}
\item Calculer la fréquence de chaque classe d'âge.
\item Donner une estimation de l'étendue des âges.
\item Calculer la vitesse moyenne en km/h d'un participant.
\item Si le nombre de participants augmente de $15\%$ l'année suivante, calculer le nouvel effectif total.
\item Rédiger une synthèse mathématique des résultats obtenus.
\end{enumerate}
\end{exobox}

\newpage
\section{Corrigé complet}
\setcounter{corr}{0}

\begin{corrbox}{Effectifs et fréquences dans une classe}
\begin{enumerate}
\item $7+12+6+3=28$, donc l'effectif total est bien $28$.
\item La fréquence des élèves venant en bus vaut $\dfrac{12}{28}=\dfrac{3}{7}\approx0{,}429$, soit environ $42{,}9\%$.
\item À pied ou à vélo : $7+3=10$ élèves, donc la fréquence vaut $\dfrac{10}{28}=\dfrac{5}{14}\approx0{,}357$.
\item Les fréquences en pourcentage sont environ $42{,}9\%$ et $35{,}7\%$.
\end{enumerate}
\end{corrbox}

\begin{corrbox}{Étendue d'une série brute}
\begin{enumerate}
\item La valeur minimale est $16$.
\item La valeur maximale est $25$.
\item L'étendue vaut $25-16=9$.
\item Les températures maximales observées varient sur un intervalle de $9^\circ$C.
\end{enumerate}
\end{corrbox}

\begin{corrbox}{Fréquences dans un tableau}
\begin{enumerate}
\item Bleu : $\dfrac{18}{60}=0{,}30=30\%$.
\item Noir : $\dfrac{15}{60}=0{,}25=25\%$.
\item Rouge ou vert : $\dfrac{12+9}{60}=\dfrac{21}{60}=0{,}35=35\%$.
\item Somme : $\dfrac{18+12+9+15+6}{60}=\dfrac{60}{60}=1$.
\end{enumerate}
\end{corrbox}

\begin{corrbox}{Histogramme à classes de même amplitude}
\begin{enumerate}
\item $8+22+34+26+10=100$.
\item L'amplitude commune est $10$ minutes.
\item Les données sont regroupées en intervalles de même amplitude : l'histogramme est donc adapté.
\item Entre $20$ et $40$ minutes : $34+26=60$ élèves, donc $\dfrac{60}{100}=60\%$.
\end{enumerate}
\end{corrbox}

\begin{corrbox}{Diagramme circulaire}
\begin{enumerate}
\item Fréquences : football $\frac{32}{80}=0{,}4$, basket $0{,}25$, tennis $0{,}20$, athlétisme $0{,}15$.
\item Angles : $0{,}4\times360=144^\circ$, $0{,}25\times360=90^\circ$, $0{,}20\times360=72^\circ$, $0{,}15\times360=54^\circ$.
\item $144+90+72+54=360$ : le diagramme est complet.
\end{enumerate}
\end{corrbox}

\begin{corrbox}{Probabilité avec un dé équilibré}
\begin{enumerate}
\item $\Omega=\{1,2,3,4,5,6\}$.
\item $\Prob(5)=\dfrac16$.
\item Les nombres pairs sont $2,4,6$, donc $\Prob=\dfrac36=\dfrac12$.
\item Les nombres supérieurs ou égaux à $4$ sont $4,5,6$, donc $\Prob=\dfrac36=\dfrac12$.
\item Ne pas obtenir $1$ : $1-\dfrac16=\dfrac56$.
\end{enumerate}
\end{corrbox}

\begin{corrbox}{Probabilité dans une urne}
\begin{enumerate}
\item Total : $4+3+5=12$ boules.
\item Rouge : $\dfrac{4}{12}=\dfrac13$.
\item Bleue ou verte : $\dfrac{3+5}{12}=\dfrac{8}{12}=\dfrac23$.
\item Ne pas tirer vert : $1-\dfrac{5}{12}=\dfrac{7}{12}$.
\end{enumerate}
\end{corrbox}

\begin{corrbox}{Deux pièces équilibrées}
\begin{enumerate}
\item Issues : $(P,P),(P,F),(F,P),(F,F)$.
\item Deux piles : $\dfrac14$.
\item Exactement une pile : deux issues, donc $\dfrac24=\dfrac12$.
\item Au moins une face : tout sauf $(P,P)$, donc $\dfrac34$.
\end{enumerate}
\end{corrbox}

\begin{corrbox}{Fonction linéaire et vitesse constante}
\begin{enumerate}
\item La distance est $d(x)=5x$.
\item $d(8)=5\times8=40$ m.
\item $1$ minute vaut $60$ s, donc $d(60)=300$ m.
\item $5x=250$, donc $x=50$ s.
\end{enumerate}
\end{corrbox}

\begin{corrbox}{Pourcentages d'évolution}
\begin{enumerate}
\item Hausse de $15\%$ : coefficient $1{,}15$, prix $120\times1{,}15=138$ euros.
\item Baisse de $20\%$ : coefficient $0{,}80$, prix $120\times0{,}8=96$ euros.
\item Hausse de $8\%$ : coefficient $1{,}08$.
\item Baisse de $35\%$ : coefficient $0{,}65$.
\end{enumerate}
\end{corrbox}

\begin{corrbox}{Volume d'une boule}
\begin{enumerate}
\item $V=\dfrac43\pi r^3$.
\item Pour $r=6$, $V=\dfrac43\pi\times216=288\pi$ cm$^3$.
\item $288\pi\approx904{,}8$ cm$^3$.
\item Le rayon est en cm, donc le volume est en cm$^3$.
\end{enumerate}
\end{corrbox}

\begin{corrbox}{Vitesse et conversion}
\begin{enumerate}
\item $72\div3{,}6=20$, donc $72$ km/h $=20$ m/s.
\item En $3$ s : $20\times3=60$ m.
\item En $12$ s : $20\times12=240$ m.
\item Le temps est en secondes ; il faut donc une vitesse en m/s pour obtenir directement des mètres.
\end{enumerate}
\end{corrbox}

\begin{corrbox}{Retrouver des effectifs à partir de fréquences}
\begin{enumerate}
\item $0{,}45+0{,}25+0{,}20+0{,}10=1$.
\item Bus : $480\times0{,}45=216$ ; voiture : $120$ ; à pied : $96$ ; vélo : $48$.
\item $216+120+96+48=480$.
\item Le tableau d'effectifs est donc : bus $216$, voiture $120$, à pied $96$, vélo $48$.
\end{enumerate}
\end{corrbox}

\begin{corrbox}{Lire un diagramme en bâtons décrit par un tableau}
\begin{enumerate}
\item Au moins $40$ : $8+6+5+3=22$.
\item Au plus $42$ : $2+5+7+8+6+5=33$.
\item De $38$ à $41$ inclus : $5+7+8+6=26$.
\item Au moins $41$ : $6+5+3=14$, fréquence $\dfrac{14}{36}=\dfrac{7}{18}$.
\end{enumerate}
\end{corrbox}

\begin{corrbox}{Construire un histogramme complet}
\begin{enumerate}
\item Chaque classe a une amplitude de $5$ kg.
\item Une échelle possible : $1$ cm pour $5$ colis sur l'axe vertical.
\item Les rectangles accolés ont pour bases les classes et pour hauteurs les effectifs $6,14,18,9,3$.
\item Masse au moins $10$ kg : $18+9+3=30$, fréquence $\dfrac{30}{50}=60\%$.
\end{enumerate}
\end{corrbox}

\begin{corrbox}{Étendue à partir d'un histogramme}
\begin{enumerate}
\item Tableau : classes données et effectifs $3,9,15,10,3$.
\item Total : $3+9+15+10+3=40$.
\item Estimation de l'étendue : $190-140=50$ cm.
\item Les données exactes ne sont pas connues ; on ne connaît que les classes, donc l'étendue exacte peut être plus petite.
\end{enumerate}
\end{corrbox}

\begin{corrbox}{Deux tirages avec remise}
\begin{enumerate}
\item Le tableau contient les couples $(B,B)$, $(B,V_1)$, $(B,V_2)$, $(V_1,B)$, $(V_1,V_1)$, $(V_1,V_2)$, $(V_2,B)$, $(V_2,V_1)$, $(V_2,V_2)$.
\item Il y a $3\times3=9$ issues.
\item Deux violettes : $(V_1,V_1)$, $(V_1,V_2)$, $(V_2,V_1)$, $(V_2,V_2)$.
\item La probabilité est donc $\dfrac49$.
\end{enumerate}
\end{corrbox}

\begin{corrbox}{Deux enfants}
\begin{enumerate}
\item Issues : $(F,F),(F,G),(G,F),(G,G)$.
\item Deux garçons : $\dfrac14$.
\item Au moins une fille : $(F,F),(F,G),(G,F)$, donc $\dfrac34$.
\item L'événement contraire est « deux garçons », donc $1-\dfrac14=\dfrac34$.
\end{enumerate}
\end{corrbox}

\begin{corrbox}{Simulation d'un dé}
\begin{enumerate}
\item Fréquences : $0{,}165$, $0{,}171$, $0{,}164$, $0{,}169$, $0{,}166$, $0{,}165$ au millième près.
\item Probabilité théorique : $\dfrac16\approx0{,}167$.
\item Les fréquences sont toutes proches de $0{,}167$.
\item Sur un grand nombre de lancers, les fréquences observées se stabilisent autour des probabilités théoriques.
\end{enumerate}
\end{corrbox}

\begin{corrbox}{Coefficient multiplicateur global}
\begin{enumerate}
\item Hausse de $10\%$ : $1{,}10$.
\item Baisse de $10\%$ : $0{,}90$.
\item Coefficient global : $1{,}10\times0{,}90=0{,}99$.
\item Le prix final vaut $99\%$ du prix initial ; il n'est donc pas égal au prix initial.
\end{enumerate}
\end{corrbox}

\begin{corrbox}{Fonction linéaire et tableau}
\begin{enumerate}
\item Coefficient : $18\div6=3$.
\item Fonction : $y=3x$.
\item Pour $x=2,5,9,12$, on obtient $y=6,15,27,36$.
\item $3x=45$, donc $x=15$.
\end{enumerate}
\end{corrbox}

\begin{corrbox}{Thalès et proportionnalité}
\begin{enumerate}
\item $\dfrac{AM}{AB}=\dfrac{AN}{AC}$.
\item $\dfrac{AM}{AB}=\dfrac4{10}=0{,}4$.
\item $\dfrac{AN}{15}=0{,}4$, donc $AN=6$ cm.
\item Le triangle $AMN$ est une réduction du triangle $ABC$ de coefficient $0{,}4$.
\end{enumerate}
\end{corrbox}

\begin{corrbox}{Débit et unités}
\begin{enumerate}
\item En $7$ min : $18\times7=126$ L.
\item $126$ L $=0{,}126$ m$^3$.
\item $540\div18=30$ min.
\item L/min multiplié par min donne L ; L divisé par L/min donne min.
\end{enumerate}
\end{corrbox}

\begin{corrbox}{Assemblage de solides}
\begin{enumerate}
\item Cylindre : $V_1=\pi\times3^2\times12=108\pi$ cm$^3$.
\item Demi-boule : $V_2=\dfrac12\times\dfrac43\pi\times3^3=18\pi$ cm$^3$.
\item Total : $126\pi$ cm$^3$.
\item $126\pi\approx396$ cm$^3$.
\end{enumerate}
\end{corrbox}

\begin{corrbox}{Justifier une formule de fréquence}
\begin{enumerate}
\item Une fréquence mesure la part d'une modalité dans le total, donc $f=n/N$.
\item Comme $0\leq n\leq N$, en divisant par $N>0$, $0\leq n/N\leq1$.
\item Si les effectifs sont $n_1,\ldots,n_k$, alors $n_1+\cdots+n_k=N$, donc $\frac{n_1}{N}+\cdots+\frac{n_k}{N}=1$.
\item $f=0{,}25$ signifie que la modalité représente un quart de la série, soit $25\%$.
\end{enumerate}
\end{corrbox}

\begin{corrbox}{Justifier l'étendue}
\begin{enumerate}
\item L'étendue est $M-m$.
\item Comme $M\geq m$, on a $M-m\geq0$.
\item Si toutes les valeurs sont égales, par exemple $7,7,7$, l'étendue vaut $0$.
\item Elle mesure l'écart entre les extrêmes, mais ne renseigne pas sur la répartition intérieure des valeurs.
\end{enumerate}
\end{corrbox}

\begin{corrbox}{Justifier l'événement contraire}
\begin{enumerate}
\item $\overline A$ se réalise exactement lorsque $A$ ne se réalise pas.
\item Une issue ne peut pas à la fois réaliser $A$ et ne pas réaliser $A$.
\item Toute issue réalise soit $A$, soit son contraire.
\item Les probabilités des deux événements se complètent pour faire $1$, donc $\Prob(\overline A)=1-\Prob(A)$.
\end{enumerate}
\end{corrbox}

\begin{corrbox}{Justifier une situation de proportionnalité}
\begin{enumerate}
\item $7{,}5/3=2{,}5$, $12{,}5/5=2{,}5$, $20/8=2{,}5$, $30/12=2{,}5$.
\item Les quotients sont égaux : le tableau est proportionnel.
\item Le coefficient vaut $2{,}5$.
\item La fonction est $y=2{,}5x$.
\end{enumerate}
\end{corrbox}

\begin{corrbox}{Prouver qu'une situation n'est pas proportionnelle}
\begin{enumerate}
\item $5/2=2{,}5$, $9/4=2{,}25$, $17/8=2{,}125$.
\item Ces quotients ne sont pas égaux.
\item Il n'existe donc pas de coefficient constant : le tableau n'est pas proportionnel.
\item Deux grandeurs peuvent augmenter ensemble sans que leur quotient reste constant.
\end{enumerate}
\end{corrbox}

\begin{corrbox}{Démontrer le coefficient d'une hausse}
\begin{enumerate}
\item La quantité ajoutée est $\dfrac{p}{100}V$.
\item La valeur finale est $V+\dfrac{p}{100}V$.
\item On factorise : $V\left(1+\dfrac{p}{100}\right)$.
\item Le coefficient multiplicateur est donc $1+\dfrac{p}{100}$.
\end{enumerate}
\end{corrbox}

\begin{corrbox}{Démontrer le coefficient d'une baisse}
\begin{enumerate}
\item La quantité retirée est $\dfrac{p}{100}V$.
\item La valeur finale est $V-\dfrac{p}{100}V$.
\item On factorise : $V\left(1-\dfrac{p}{100}\right)$.
\item Le coefficient multiplicateur est donc $1-\dfrac{p}{100}$.
\end{enumerate}
\end{corrbox}

\begin{corrbox}{Contrôle d'unités dans une vitesse}
\begin{enumerate}
\item L'unité est m/s.
\item m/s multiplié par s donne m, car les secondes se simplifient.
\item $1$ km/h $=\dfrac{1000\text{ m}}{3600\text{ s}}=\dfrac1{3{,}6}$ m/s.
\item Ainsi, pour convertir une vitesse de km/h en m/s, on divise par $3{,}6$.
\end{enumerate}
\end{corrbox}

\begin{corrbox}{Histogramme et piège de lecture}
\begin{enumerate}
\item La classe modale est \classe{10}{15}, d'effectif $42$.
\item Inférieurs à $10$ min : $10+18=28$, fréquence $\frac{28}{120}=\frac{7}{30}$.
\item Au moins $15$ min : $30+20=50$, fréquence $\frac{50}{120}=\frac{5}{12}$.
\item Non : les données sont regroupées ; on ne connaît pas les valeurs individuelles exactes.
\end{enumerate}
\end{corrbox}

\begin{corrbox}{Probabilités et dénombrement non ordonné}
\begin{enumerate}
\item Il y a $4\times4=16$ couples possibles, car il y a remise.
\item Deux rouges : $2\times2=4$ issues, donc $4/16=1/4$.
\item Même couleur : rouges $4$ issues ou bleus $4$ issues, donc $8/16=1/2$.
\item Distinguer les jetons rend les issues équiprobables et évite les doubles comptages imprécis.
\end{enumerate}
\end{corrbox}

\begin{corrbox}{Comparer deux simulations}
\begin{enumerate}
\item Premier : $14/60\approx0{,}233$ ; second : $1015/6000\approx0{,}169$.
\item Théorie : $1/6\approx0{,}167$.
\item La seconde fréquence est plus proche de $1/6$.
\item Un grand nombre d'essais donne généralement une fréquence plus stabilisée.
\end{enumerate}
\end{corrbox}

\begin{corrbox}{Hausse puis baisse}
\begin{enumerate}
\item Après hausse : $200\times1{,}25=250$ euros.
\item Après baisse : $250\times0{,}80=200$ euros.
\item Coefficient global : $1{,}25\times0{,}80=1$.
\item Comme le coefficient global vaut $1$, le prix final est égal au prix initial.
\end{enumerate}
\end{corrbox}

\begin{corrbox}{Baisse puis hausse inverse ?}
\begin{enumerate}
\item Baisse de $30\%$ : coefficient $0{,}70$.
\item Si $70$ est le prix après baisse, le prix initial était $70/0{,}70=100$.
\item Il faut passer de $70$ à $100$, soit une hausse de $30$ sur $70$ : $30/70\approx42{,}86\%$.
\item Le pourcentage est calculé sur la nouvelle base $70$, pas sur l'ancienne base $100$.
\end{enumerate}
\end{corrbox}

\begin{corrbox}{Thalès avec retour à la longueur initiale}
\begin{enumerate}
\item $AB=AM+MB=6+9=15$ cm.
\item $AM/AB=6/15=2/5$.
\item $AN/AC=2/5$, donc $8/AC=2/5$, d'où $AC=20$ cm.
\item $NC=AC-AN=20-8=12$ cm.
\end{enumerate}
\end{corrbox}

\begin{corrbox}{Boîte cylindrique et boules}
\begin{enumerate}
\item Boîte : $V_c=\pi\times5^2\times30=750\pi$ cm$^3$.
\item Trois boules : $3\times\frac43\pi\times5^3=500\pi$ cm$^3$.
\item Volume restant : $750\pi-500\pi=250\pi$ cm$^3$.
\item $250\pi\approx785$ cm$^3$ ; le résultat est positif et inférieur au volume de la boîte, il est cohérent.
\end{enumerate}
\end{corrbox}

\begin{corrbox}{Débit important}
\begin{enumerate}
\item $3$ min $=180$ s.
\item Volume : $328\times180=59040$ m$^3$.
\item En litres : $59040\times1000=59\,040\,000$ L.
\item $59\,040\,000=5{,}904\times10^7$ L.
\end{enumerate}
\end{corrbox}

\begin{corrbox}{Choisir la bonne méthode}
\begin{enumerate}
\item $300\times0{,}42=126$ élèves.
\item $\frac27\times126=36$ élèves.
\item Parmi tous les élèves : $36/300=0{,}12=12\%$.
\item Les calculs de pourcentage et de fraction d'une quantité relèvent de la proportionnalité.
\end{enumerate}
\end{corrbox}

\begin{corrbox}{Données regroupées et incertitude}
\begin{enumerate}
\item $70/200=0{,}35=35\%$.
\item Estimation : $100-0=100$ minutes.
\item Non, car les données sont regroupées par classes.
\item Oui : les classes à partir de $40$ minutes sont \classe{40}{60}, \classe{60}{80}, \classe{80}{100}, donc $70+40+15=125$ utilisateurs.
\end{enumerate}
\end{corrbox}

\begin{corrbox}{Enquête statistique complète}
\begin{enumerate}
\item $100+200+400+1100+700=2500$.
\item Fréquences : $4\%$, $8\%$, $16\%$, $44\%$, $28\%$.
\item Les rectangles accolés ont même largeur $4$ ans et hauteurs $100,200,400,1100,700$.
\item La classe la plus représentée est \classe{30}{34}.
\item Estimation de l'étendue : $38-18=20$ ans.
\item La majorité des personnes répondantes ont trouvé un emploi correspondant à leur qualification entre $30$ et $38$ ans. La classe \classe{30}{34} est dominante. Les données étant regroupées, les âges exacts ne sont pas connus.
\end{enumerate}
\end{corrbox}

\begin{corrbox}{Jeu à deux épreuves}
\begin{enumerate}
\item Les issues sont les $9$ couples formés avec $B,V_1,V_2$.
\item Deux violettes : $4$ issues sur $9$, donc $4/9$.
\item Aucune violette signifie $(B,B)$, donc $1/9$.
\item Gagner signifie au moins une violette : $1-1/9=8/9$.
\item Les issues gagnantes sont toutes sauf $(B,B)$ : il y en a $8$, donc $8/9$.
\end{enumerate}
\end{corrbox}

\begin{corrbox}{Trajet, vitesse et temps de réaction}
\begin{enumerate}
\item $80/3{,}6\approx22{,}2$ m/s.
\item Distance de réaction : $22{,}2\times1\approx22{,}2$ m.
\item Distance totale : $22{,}2+38\approx60{,}2$ m.
\item $90/3{,}6=25$ m/s, donc distance de réaction $25$ m.
\item À $90$ km/h, la distance de réaction est plus grande ; la distance totale serait donc plus grande si la distance de freinage restait identique.
\end{enumerate}
\end{corrbox}

\begin{corrbox}{Réservoir composé}
\begin{enumerate}
\item Cylindre : $\pi\times4^2\times10=160\pi$ dm$^3$.
\item Demi-boule : $\frac12\times\frac43\pi\times4^3=\frac{128}{3}\pi$ dm$^3$.
\item Total : $160\pi+\frac{128}{3}\pi=\frac{608}{3}\pi$ dm$^3$.
\item Comme $1$ dm$^3=1$ L, le volume est $\frac{608}{3}\pi$ L.
\item Valeur approchée : environ $637$ L.
\end{enumerate}
\end{corrbox}

\begin{corrbox}{Sondage, fréquences et évolution}
\begin{enumerate}
\item $600\times0{,}35=210$ adhérents.
\item $600\times1{,}12=672$ adhérents.
\item $672\times0{,}40=268{,}8$. Dans un contexte réel, cela correspondrait à environ $269$ adhérents ; les données en pourcentage peuvent être arrondies.
\item Augmentation approximative : $269-210=59$ adhérents.
\item Le total augmente et la proportion de jeunes augmente aussi ; l'effectif des moins de $18$ ans augmente fortement.
\end{enumerate}
\end{corrbox}

\begin{corrbox}{Fonction linéaire et géométrie}
\begin{enumerate}
\item $L'=1{,}8L$.
\item $4\mapsto7{,}2$ cm ; $7{,}5\mapsto13{,}5$ cm ; $12\mapsto21{,}6$ cm.
\item $1{,}8L=45$, donc $L=25$ cm.
\item Toutes les longueurs sont multipliées par le même coefficient $1{,}8$ : elles sont proportionnelles.
\item Les volumes augmenteraient plus vite que les longueurs ; il ne suffit pas de les multiplier par $1{,}8$.
\end{enumerate}
\end{corrbox}

\begin{corrbox}{Interprétation d'une expérience simulée}
\begin{enumerate}
\item Fréquences : $14/20=0{,}70$, $107/200=0{,}535$, $10084/20000=0{,}5042$.
\item Probabilité théorique de pile : $1/2=0{,}5$.
\item Les fréquences se rapprochent globalement de $0{,}5$ quand le nombre de lancers augmente.
\item La troisième simulation repose sur $20000$ essais ; elle est donc plus stabilisée.
\item Les fréquences observées peuvent varier sur de petits échantillons, mais elles tendent à se stabiliser autour de la probabilité théorique sur un grand nombre d'essais.
\end{enumerate}
\end{corrbox}

\begin{corrbox}{Problème bilan transversal}
\begin{enumerate}
\item Fréquences : $120/800=15\%$, $260/800=32{,}5\%$, $220/800=27{,}5\%$, $140/800=17{,}5\%$, $60/800=7{,}5\%$.
\item Estimation de l'étendue : $60-10=50$ ans.
\item $36$ min $=0{,}6$ h, donc $v=5/0{,}6=8{,}33$ km/h environ.
\item Nouvel effectif : $800\times1{,}15=920$ participants.
\item La course rassemble surtout les $20$-$40$ ans, avec une vitesse moyenne d'environ $8{,}33$ km/h. Une hausse de $15\%$ porterait l'effectif à $920$ participants.
\end{enumerate}
\end{corrbox}

\end{document}